%============================================================ 4. 제안 방법 (압축판)
%  원판(sections_real/04_method.tex)에서 구현 세부 수식을 산문으로 내리고, 결과에 쓰이지
%  않은 반사실 크레딧 변형을 덜어냈다. 남긴 식은 방법의 정체성을 정하는 것들이다.
\section{제안 방법}
\label{sec:method}

\subsection{2계층 비동기 구조}
\label{sec:arch}

제안 체계는 방어 판단을 전략 계층과 기동 계층으로 분리한다(Fig.~\ref{fig:architecture}).
전략 계층은 어느 방어정이 어느 적 무리를 맡는가를 정하고, 기동 계층은 배정된
무리를 막기 위해 해면 어디에 그물벽을 놓는가를 정한다. 두 계층은 같은 주기로 결정을
내리지만 응답 속도가 다르다. 기동 계층의 순전파는 순간에 끝나는 반면, 전략 계층의
언어모델 호출은 응답에 지연이 발생한다(\ref{sec:res_latency}절). 느린 전략 계층 호출만
비동기로 분리하여, 호출이 진행되는 동안 기동 계층은 직전 배정으로 결정을 계속하고 응답이
도착하면 다음 결정부터 새 배정을 반영한다.

\subsection{적 무리 군집화}
\label{sec:cluster}
\begin{figure}[!htbp]
\centering
\includegraphics[width=\columnwidth]{fig5.png}
\caption{방위각 정렬, 군집화, 군집 배정}
\label{fig:clustering}
\end{figure}

두 판단 계층은 개별 적선이 아니라 무리를 단위로 배정하므로, 매 결정마다 생존 적선을 모선
기준 방위 $\vartheta_m$ 으로 묶어 최대 $C$ 개의 무리로 요약한다(Fig.~\ref{fig:clustering}).
방위 구간을 고정해 나누면 무리가 경계에 걸려 쪼개지므로, 정렬한 방위의 이웃 간격 $g_{(q)}$ 로
무리 수를 데이터에 맞춘다. 방위는 순환값이라 간격도 순환으로 센다. 임계
$\delta_{\mathrm{gap}}$ 을 넘는 간격의 개수가 활성 무리 수 $\hat{C}$ 이며,
$1$ 이상 $\min(C, n)$ 이하로 자른다.
\begin{equation}
\hat{C} = \mathrm{clip}\!\Big(\big|\{\, q : g_{(q)} > \delta_{\mathrm{gap}} \,\}\big|,\; 1,\; \min(C, n)\Big)
\label{eq:ncluster}
\end{equation}
가장 큰 간격 $\hat{C}$ 개를 경계로 삼아 방위 순으로 라벨을 붙여 무리별 구성원을 정하고, 각
무리를 척수 $n_j$, 중심 $\mathbf{c}_j$, 각도 퍼짐 $\Delta\vartheta_j$, 모선 방향 접근 속도
$v^{\mathrm{app}}_j$ 의 네 양으로 요약한다.
% ★ Fig.6(fig:observation) 은 2단 figure* 라 선언된 쪽에 놓일 수 없어 여기(§4.2 끝)에서
%   선언해야 §4.4.1 이 시작하는 쪽 상단에 앉는다.
\begin{figure*}[!tp]
\centering
\includegraphics[width=0.62\textwidth]{fig6.png}
\caption{래스터화된 15채널 관측의 예(파상, 통영 매물도 북측)}
\label{fig:observation}
\end{figure*}

이 규칙에서 집중은 무리 하나, 양동은 셋, 파상은 대개 단마다 하나로 요약된다. 척수와 중심은
\ref{sec:heuristic}절의 위협도와 요격점을, 네 요약량 전부는 \ref{sec:commander}절의
지휘관 입력을 구성한다.

\subsection{기하 휴리스틱 기준 방법}
\label{sec:heuristic}

무리 요약 위에서 같은 문제를 기하 규칙만으로 푸는 기준 방법으로, \ref{sec:results}절의 비교
기준이자 \ref{sec:bc}절의 행동 복제 목표다. 무리 $j$ 의 위협도 $\tau_j$ 는 척수 $n_j$ 에 모선
근접도를 곱한 양이고, 요격점 $\mathbf{I}_j$ 는 무리 중심에서 모선을 향해 선분의 $\lambda_{\mathrm{I}}$
비율만큼 옮긴 점이다.
\begin{equation}
\begin{aligned}
\tau_j &= n_j \cdot \mathrm{clip}\!\left(1 - \frac{\lVert \mathbf{c}_j - \mathbf{m}\rVert}{R_w},
         \, 0.05,\, 1\right),\\
\mathbf{I}_j &= \mathbf{c}_j + \lambda_{\mathrm{I}}\,(\mathbf{m} - \mathbf{c}_j), \qquad \lambda_{\mathrm{I}} = 0.55
\end{aligned}
\label{eq:threat}
\end{equation}
하한은 먼 무리도 척수에 비례하는 최소 위협을 갖게 한다. 방어정이 적보다 느려 무리 중심
가까이 잡은 요격점에는 적보다 늦게 도달하므로, 요격점은 모선 쪽으로 치우쳐 둔다.

위협 상위 $P$ 개 무리에 대해 탐욕적 1:1 배정을 한다. 비용은 방어정 위치 $\mathbf{a}_p$ 에서
요격점까지의 거리 $d_{pj}$ 에서, 직전 결정의 담당 $j_p^{\text{prev}}$ 를 유지하는 쌍에만
보너스 $\beta$ 를 감한 값이다. $\beta$ 는 해역 규모에 맞먹는 크기로 두어, 거리 차이가
그만큼 벌어지기 전에는 직전 배정이 유지되게 한다.
\begin{equation}
\tilde{d}_{pj} = \begin{cases}
  d_{pj} - \beta, & j = j_p^{\text{prev}}\\
  d_{pj},          & \text{그 밖}
\end{cases}
\label{eq:sticky}
\end{equation}
활성 무리가 $P$ 보다 적으면 남는 방어정은 부하가 가장 큰 무리에 안쪽 층으로 덧붙인다. 층은 반경 증분 $\lambda_{\mathrm{L}}$ 과 방위차만큼 어긋나게 둔다. 그물벽은 요격점을 가운데 둔, 접근 방향에 수직인
길이 $L_{\mathrm{net}}$ 의 선분이므로 경유점은 그 양 끝 $K = 2$ 개이며, 학습 정책과 같은
격자·유효 마스크(\ref{sec:validmask}절) 위의 최근접 픽셀로 스냅한다.

따라서 두 방법의 행동공간은 같고 다른 것은 결정 규칙뿐이다.

\subsection{점수맵 기반 기동 정책}
\label{sec:policy}

\subsubsection{다채널 격자 관측}
\label{sec:obs}

기동 계층의 입력은 좌표 목록이 아니라, 모선을 중앙에 두고 관측 반폭 $E$ 의 전장을
$n_{\mathrm{g}} \times n_{\mathrm{g}}$ 격자로 나눈 다채널 지도이다. 총 15개 채널은 모든
방어정이 공유하는 전역 9채널, 방어정마다 다른 3채널, CoordConv~\citep{liu2018coordconv}
좌표 3채널로 구성되며(Fig.~\ref{fig:observation}), 각 채널의 정의와 값 범위를
Table~\ref{tab:channels}에 정리하였다.

\begin{table*}[!tp]
\centering
\caption{15채널 격자 관측의 구성}
\label{tab:channels}
\footnotesize
\renewcommand{\arraystretch}{1.15}
\begin{tabular}{@{}clllc@{}}
\toprule
\# & 묶음 & 채널 & 정의 & 범위 \\
\midrule
1  & \multirow{9}{*}{전역}   & 적 밀도        & 픽셀 내 생존 적선 수$/M$                                & $[0,1]$ \\
2  &                         & 적 진행 $x$    & 픽셀 내 적선 헤딩 단위벡터의 평균                      & $[-1,1]$ \\
3  &                         & 적 진행 $y$    & 위와 같음                                              & $[-1,1]$ \\
4  &                         & 적 위협도      & 모선 도달 임박도의 픽셀 내 최대                        & $[0,1]$ \\
5  &                         & 아군 위치      & 픽셀 내 생존 방어정 수$/P$                             & $[0,1]$ \\
6  &                         & 설치된 그물    & 물리 격자의 그물 점유를 최대값 풀링                     & $\{0,1\}$ \\
7  &                         & 원점           & 모선 반경 안의 픽셀                                    & $\{0,1\}$ \\
8  &                         & 요격 구역      & $r_{\min} \le \lVert \mathbf{x}_i - \mathbf{m} \rVert \le r_{\max}$ 인 픽셀 & $\{0,1\}$ \\
9  &                         & 육지           & 고해상도 육지 마스크의 면적비  & $[0,1]$ \\
\midrule
10 & \multirow{3}{*}{방어정별} & 자기 위치     & 방어정 $p$ 가 놓인 픽셀                                & $\{0,1\}$ \\
11 &                         & 배정 요격점    & $\mathbf{I}_p$ 가 놓인 픽셀 (미배정이면 전부 0)        & $\{0,1\}$ \\
12 &                         & 유효 영역      & 유효 마스크 $\mathcal{V}_p$, 식~\eqref{eq:validset}    & $\{0,1\}$ \\
\midrule
13 & \multirow{3}{*}{좌표}   & $x$            & $(x_i - m_x)/E$                                        & $[-1,1]$ \\
14 &                         & $y$            & $(y_i - m_y)/E$                                        & $[-1,1]$ \\
15 &                         & $r$            & $\mathrm{clip}(\lVert \mathbf{x}_i - \mathbf{m} \rVert / E, 0, 1)$ & $[0,1]$ \\
\bottomrule
\end{tabular}
\end{table*}

픽셀 크기는 $\ell = 2E/n_{\mathrm{g}}$ 이고, 월드 좌표
$\mathbf{x} = (x, y)$ 는 모선 위치 $\mathbf{m} = (m_x, m_y)$ 를 중심으로 다음 픽셀 인덱스
$i = (\iota_x, \iota_y)$ 로 사상된다.
\begin{equation}
\begin{aligned}
\iota_x &= \mathrm{clip}\!\left(\left\lfloor \frac{x - m_x + E}{\ell} \right\rfloor,\, 0,\, n_{\mathrm{g}}-1\right),\\
\iota_y &= \mathrm{clip}\!\left(\left\lfloor \frac{y - m_y + E}{\ell} \right\rfloor,\, 0,\, n_{\mathrm{g}}-1\right)
\end{aligned}
\label{eq:world2pix}
\end{equation}
적 관련 채널은 각 개체의 연속 위치를 식~\eqref{eq:world2pix}로 픽셀에 사상한 뒤 같은 픽셀에
속한 개체를 집계하여 값을 만든다. 밀도는 픽셀 내 적선 수를 전체 수 $M$ 으로, 진행 방향은
헤딩 단위벡터의 픽셀 내 평균으로, 위협도는 적선이 에피소드 상한 안에 달릴 수 있는 최대
거리로 모선까지 거리를 정규화해 1에서 뺀 값의 픽셀 내 최대로 둔다. 아군 위치 채널은 같은
방식으로 $P$ 로 나누고, 설치된 그물 채널은 고해상도 물리 격자의 최대값 풀링, 육지 채널은
OpenStreetMap 해안선·섬 폴리곤을 래스터화한 실해역 마스크의 면적비 다운샘플이고, 충돌 판정은
다운샘플 전 물리 격자에서 한다. 좌표 채널은 픽셀 중심의 모선 상대 위치와 거리를 반폭 $E$ 로
나눈 값이다. 격자에 담지 않는 방어정 $p$ 의 자체 상태 벡터 $\mathbf{o}_p \in \mathbb{R}^{8}$
는 같은 반폭으로 정규화한 자기 위치와 배정 요격점, 헤딩 단위벡터, 잔여 그물 비율, 전개 중
여부로 구성된다. 합성곱은 평행이동 등변이라 절대 위치가 흐려지므로 좌표 채널을 함께 넣어
위치를 보존한다. 모든 입력은 길이·개수·시간으로 나뉜 무차원 양이고 격자 지도는 개체 수와
순서에 불변이므로, 해역 크기·격자 해상도·선박 수가 바뀌어도 입출력의 형태가 같다.

\subsubsection{FiLM 조건화 U-Net과 점수맵}

\begin{figure*}[!tbp]
\centering
\includegraphics[width=0.66\textwidth]{fig8.png}
\caption{점수맵 정책의 구조: 질의·특징맵 내적과 칸 안 오프셋}
\label{fig:policy}
\end{figure*}

관측은 U-Net 구조~\citep{ronneberger2015unet}를 지나 입력과 같은 해상도의 특징 지도로
변환된다(Fig.~\ref{fig:policy}). 전역 평균 풀링으로 압축한 전황 벡터는
FiLM~\citep{perez2018film} 변조로 특징에 실려, 국소 합성곱이 옮기지 못하는 전황을 모든 위치에
전달한다.

각 방어정 $p$ 는 자신의 상태 벡터 $\mathbf{o}_p$ 로부터 MLP 로 질의 $\mathbf{q}_p$ 를 만들고,
아핀 사영 $\mathbf{W}_q$, $\mathbf{b}_q$ 를 거쳐 픽셀 $i$ 의 특징 $\mathbf{f}_i$ 와 내적하여
점수 $s_{p,i}$ 를 얻는다.
\begin{equation}
s_{p,i} = \frac{\langle \mathbf{W}_q \mathbf{q}_p + \mathbf{b}_q,\, \mathbf{f}_i \rangle}{\sqrt{d}}, \qquad d = 32
\label{eq:score}
\end{equation}
어텐션~\citep{vaswani2017attention}과 같은 형식의 점수로 격자 위의 원소를 지목하는 포인터
구조다~\citep{vinyals2015pointer}. 관측 텐서는 방어정을 묶어 결정마다 한 번의 순전파로 통과하며, 가중치는 모두가 공유하므로 모델 크기는 방어정 수에 불변이다. 합성곱 블록의 그룹
정규화~\citep{wu2018groupnorm}는 배치 통계를 쓰지 않아 묶음 크기와 무관하게 같은 연산이다.

\subsubsection{유효 마스크와 픽셀 지목}
\label{sec:validmask}

점수맵 전체가 후보가 되는 것은 아니다. 육지 픽셀 집합을 $\mathcal{L}_{\mathrm{land}}$, 모선에서
본 요격점 방위와 픽셀 방위의 차이를 $[0^\circ, 180^\circ]$ 로 접은 값을
$|\vartheta_i - \vartheta_p|_{\circ}$ 라 하면, 방어정 $p$ 의 유효 픽셀 집합 $\mathcal{V}_p$ 는
환상 요격 구역, 육지 밖, 요격점 반경, 무리 방위라는 네 조건의 교집합이다(설정값은
Table~\ref{tab:validmask_param}).
\begin{equation}
\begin{aligned}
\mathcal{V}_p = \big\{\, i :\; & r_{\min} \le \lVert \mathbf{x}_i - \mathbf{m} \rVert \le r_{\max},\;
  i \notin \mathcal{L}_{\mathrm{land}},\\
  & \lVert \mathbf{x}_i - \mathbf{I}_p \rVert \le r_{\mathrm{g}},\;
  |\vartheta_i - \vartheta_p|_{\circ} \le \vartheta_{\mathrm{g}} \,\big\}
\end{aligned}
\label{eq:validset}
\end{equation}
마스크는 로짓 단계에서 적용한다. 유효 집합 밖의 점수를 $-\infty$ 로 바꾼 마스크 점수
$\tilde{s}_{p,i}$ 를 소프트맥스에 넣어 픽셀 선택 확률을 얻는다.
\begin{equationsmall}
\tilde{s}_{p,i} = \begin{cases} s_{p,i}, & i \in \mathcal{V}_p \\ -\infty, & i \notin \mathcal{V}_p, \end{cases}
\quad
\pi_\theta(i \mid \mathbf{s}, p) = \frac{e^{\tilde{s}_{p,i}}}{\sum_{j \in \mathcal{V}_p} e^{\tilde{s}_{p,j}}}
\label{eq:masked_softmax}
\end{equationsmall}
안전·임무 제약이 손실 항이 아니라 행동 분포에 직접 부과되므로 위반은 애초에 표현 불가능하다.
학습된 체크포인트에서 격자 대부분은 갈 수 없는 자리로 남으며(Fig.~\ref{fig:validmask}), 연속
좌표 회귀에는 이 구조를 표현할 수단이 없다. 범주형 표본은 유효 영역을 벗어나지 않고 엔트로피가
유한하므로 \ref{sec:training}절의 그룹 후보들이 실제로 서로 다른 칸을 밟는다.

%============================================================ 유효 마스크·픽셀 지목 설정값 (국방발표 판)
%  §4.4.3 본문에 흩어져 있던 수치를 모은 표. 길이는 체계가 스케일 불변이므로 관측
%  픽셀 한 변 ell 기준의 비율로 적는다(기준 인스턴스의 m 값은 시나리오 표). 본문에는 기호와 논지만 남기고 값은 전부 여기에 둔다.
%  값 출처: boatattack_sim/env/config.py (cell_r_min/r_max, cnn_gate_r, cnn_gate_bear,
%           cnn_valid_min, cnn_nets, cnn_dup_r, net_max_len), 유효 비율은 학습 체크포인트 실측.
\begin{table}[htbp]
\centering
\caption{유효 마스크와 픽셀 지목의 설정값(\ref{sec:validmask}절)}
\label{tab:validmask_param}
\footnotesize
\resizebox{\ifdim\width>\linewidth\linewidth\else\width\fi}{!}{%
\begin{tabular}{llr}
\toprule
구분 & 기호 / 항목 & 값 \\
\midrule
\multirow{4}{*}{유효 픽셀 조건}
 & 요격 환형 안·바깥 반경 $r_{\min}$, $r_{\max}$ & $1.6\,\ell$ / $18\,\ell$ \\
 & 요격점 반경 조건 $r_{\mathrm{g}}$             & $12\,\ell$ \\
 & 무리 방위 조건 $\vartheta_{\mathrm{g}}$        & $60^\circ$ \\
 & 유효 픽셀 수 하한                              & $150$ 칸 \\
\midrule
\multirow{3}{*}{픽셀 지목}
 & 방어정 1척당 경유점 수 $K$                     & $2$ \\
 & 두 점 간격 하한 $r_{\mathrm{dup}}$             & $0.8\,\ell$ \\
 & 두 점 간격 상한                                & $1.5\,L_{\mathrm{net}}$ \\
\midrule
실측
 & 결정당 평균 유효 비율                          & $10.5\,\%$ \\
\bottomrule
\end{tabular}
}
\end{table}


\begin{figure*}[tbp]
\centering
\includegraphics[width=\textwidth]{fig9.png}
\caption{유효 마스크와 픽셀 지목(학습된 체크포인트의 실제 출력)}
\label{fig:validmask}
\end{figure*}

각 방어정은 유효 영역 안에서 점수 상위 픽셀 $K$ 개를 순차적으로 고른다. 먼저 고른 픽셀의
특징을 질의에 더해 두 번째를 고르므로 두 점이 서로를 알고 뽑히며, 간격의 상·하한이 벽이 서지
않을 만큼 가까운 쌍과 그물 길이를 넘는 벽을 함께 배제한다. 칸 안 위치는 선택 픽셀의 특징으로부터
MLP 가 내는 평균과 학습되는 표준편차의 가우시안에서 표집해 칸의 절반 폭 $\ell/2$ 안에서
정하므로, 목표점 $\mathbf{t}_k$ 의 로그확률과 엔트로피는 범주형과 가우시안 두 분포의 합이다.

선택된 두 좌표를 잇는 선분이 그물벽 경로가 되며, 방어정은 가까운 점까지 이동한 뒤 다음
구간에서 그물을 전개한다. 그물이 없거나 배정되지 않은 방어정은 정지시킨다.

\subsection{가치망 없는 다중 에이전트 그룹 상대 최적화}
\label{sec:training}

\begin{figure}[!tbp]
\centering
\includegraphics[width=\columnwidth]{fig11.png}
\caption{그룹 상대 최적화의 한 갱신(도식은 $G=4$, $H=4$)}
\label{fig:grpocredit}
\end{figure}

\subsubsection{그룹 상대 이득}

학습은 가치함수를 따로 두지 않는 그룹 상대 정책 최적화~\citep{shao2024deepseekmath}를 다중
방어정 상황에 맞게 변형하여 수행한다(Fig.~\ref{fig:grpocredit}). 매 결정 시점에서 정책으로부터
$G$개의 결합 행동 후보를 표집하되, 후보 가운데 하나($k=0$)는 휴리스틱(\ref{sec:heuristic}절)이
내는 행동으로 고정하여 기준선으로 삼는다. 각 후보를 시뮬레이터에서 공통 구간
$H = 180\,\mathrm{step}$ 만큼 실제로 진행시켜 창 보상 $R_k$ 를 얻은 뒤(\ref{sec:reward}절),
휴리스틱 후보의 보상 $R_0$ 를 기준으로 상대 비교하여 이득 $A_k$ 를 산정한다. 그룹 표준편차로 나누되 $\epsilon$ 을 더해 0 나눗셈을 막는다.
\begin{equation}
A_k = \frac{R_k - R_0}{\mathrm{std}(\{R_j\}_{j=0}^{G-1}) + \epsilon}
\label{eq:adv}
\end{equation}
휴리스틱을 넘어서는 후보만 양의 이득을 받고, 그룹 표준편차로 나누므로 이득은 크기가 없는
순위 정보만 남는다. 어떤 선택이 더 나았는지를 실제 결과로 직접 비교하므로 가치망과 그 추정
편향·학습 불안정이 없다.

한 갱신은 $N$개의 병렬 전장(월드)에서 동시에 수행된다. 후보 보상이 모두 같은 월드는 이득이
정의되지 않으므로 유효 월드 마스크 $w_n$ 으로 갱신에서 제외한다. 월드 $n$ 의 상태를
$\mathbf{s}_n$, 후보 $k$ 의 이득을 $A_{n,k}$, 그 후보에서 방어정 $p$ 의 행동을
$a^{(k)}_{n,p}$ 라 하면, 정책경사 손실은 유효 월드의 모든 후보·모든 방어정에 대해 이득으로
가중한 로그확률의 음수이다.
\begin{equationfn}
\mathcal{L}_{\mathrm{PG}}(\theta) = -\frac{\sum_{n} w_n \sum_{k=0}^{G-1} A_{n,k} \sum_{p=1}^{P} \log \pi_\theta\big(a^{(k)}_{n,p} \mid \mathbf{s}_n\big)}{G \sum_n w_n}
\label{eq:lpg}
\end{equationfn}
여기에 엔트로피 정칙화를 더한 총손실을 최소화하며, 기울기 노름을 클립한 뒤 갱신한다.

\label{sec:bc}
\label{sec:schedule}
그룹 상대 이득은 후보들이 갈릴 때만 신호를 주므로, 학습에 앞서 휴리스틱
행동(\ref{sec:heuristic}절)을 목표로 짧은 행동 복제로 정책을 부팅한다. 목표는 원-핫이 아니라
휴리스틱 픽셀을 중심으로 한 폭 좁은 가우시안을 유효 픽셀 위에서 정규화한 가중이다. 옵티마이저는 Adam 이며, 학습률은 $10^{-4}$ 에서 그 $5\,\%$ 까지
단일 주기 코사인 어닐링~\citep{loshchilov2017sgdr}으로, 엔트로피 계수는 $0.01$ 에서 그
$10\,\%$ 까지 선형으로, 총 800 업데이트의 같은 예산 축 위에서 함께 줄인다. 나머지 상수는
Table~\ref{tab:hparams}에 모았다.

\subsubsection{보상 설계}
\label{sec:reward}

한 후보의 창 보상은 결정 시점부터 $H$\,step 동안 일어난 일을 하나의 스칼라로 합산한 팀 공유
보상이며, 이벤트 항 $R_\mathrm{event}$, 창의 시작·끝 상태에서 잰 퍼텐셜 $\Phi$ 의
할인($\gamma$) 차분, 조밀 항 $R_\mathrm{dense}$ 로 구성된다. 세 항의 계수는
Table~\ref{tab:reward_coef}에 모았다.
\begin{equation}
R = R_\mathrm{event} + \big[\gamma\,\Phi(\mathbf{s}_{t+H}) - \Phi(\mathbf{s}_t)\big] + R_\mathrm{dense}
\label{eq:reward}
\end{equation}

\begin{table}[tp]
\centering
\caption{보상 계수}
\label{tab:reward_coef}
\footnotesize
\renewcommand{\arraystretch}{1.15}
\begin{tabular}{@{}llr@{\hspace{2mm}}l@{}}
\toprule
항 & 기호 & 값 & 의미 \\
\midrule
\multirow{6}{*}{이벤트}
 & $r_{\mathrm{cap}}$  & $+13.5$ & 포획 1건 \\
 & $r_{\mathrm{brc}}$  & $-8.0$  & 돌파 1건 \\
 & $r_{\mathrm{col}}$  & $-30$   & 충돌·육지 충돌 1건 \\
 & $r_{\mathrm{tng}}$  & $-13$   & 그물 접촉 1건 \\
 & $r_{\mathrm{wipe}}$ & $+8$    & 창 안 종료 시 전멸 \\
 & $r_{\mathrm{surv}}$ & $-5$    & 창 안 종료 시 잔존 적 1척당 \\
\midrule
\multirow{3}{*}{퍼텐셜}
 & $w_{\mathrm{threat}}$ & $3.0$  & 모선 근접도 합의 가중 \\
 & $R_w$                 & 25\,$\ell$ & 근접도 정규화(교전 해역 반폭) \\
 & $\gamma$              & $0.99$ & 할인 \\
\midrule
\multirow{10}{*}{조밀}
 & $w_{\mathrm{path}}$ & $0.5$  & 이동거리 비용 \\
 & $w_{\mathrm{net}}$  & $0.3$  & 그물 소모 비용 \\
 & $w_{\mathrm{turn}}$ & $0.3$  & 선회량 비용 \\
 & $w_{\mathrm{cons}}$ & $0.3$  & 목표점 일관성 비용 \\
 & $w_{\mathrm{prox}}$ & $0.75$ & 모선·타 방어정 근접 비용 \\
 & $w_{\mathrm{cov}}$  & $0.1$  & 커버리지 보너스 \\
 & $w_{\mathrm{ray}}$  & $1.6$  & 배율판 보너스 \\
 & $R_{\mathrm{m}}$      & 3.2\,$\ell$ & 모선 안전 반경 \\
 & $\varrho_{\mathrm{m}}$ & 2.0\,$\ell$ & 모선 근접 영향 반경 \\
 & $\varrho_{\mathrm{a}}$ & 2.4\,$\ell$ & 타 방어정 근접 영향 반경 \\
\bottomrule
\end{tabular}
\end{table}

이벤트 항은 창 안의 포획·돌파·충돌·그물 접촉 건수에 고정 가중을 곱한 합이며, 창 안에서
에피소드가 끝나면 전멸 보너스와 잔존 적 1척당 벌점을 더한다. 퍼텐셜 $\Phi$ 는 생존 적선의
모선 근접도 합에 가중을 곱해 부호를 바꾼 것이라 적이 포획되거나 멀어지면 오르며, 할인 차분 형태로 주어지므로 최적 정책을
보존한다~\citep{ng1999shaping}. 조밀 항은 기동 품질에 대한 비용·보너스의 합으로, 이동거리·소모
그물 수·선회량·직전 결정 대비 목표점 이동과 모선·타 방어정 근접도를 벌하고, 설치된 그물이 적의
접근로를 막는 비율(커버리지)을 보상한다. 포획은 평가 창의 $46.8\,\%$ 에서만 일어나므로,
나머지 창에서는 퍼텐셜·커버리지 항이 후보 순위를 가르는 주된 신호다.

\subsection{언어모델 기반 전략 계층}
\label{sec:commander}
\label{sec:schema}

전략 계층의 입력은 전장 상태·임무 목표·교전 규칙·자산 상태이고, 출력은 무리별 배정과 대기(HOLD)이며, 그 배정이 기동 계층으로 전달되어 각
방어정의 그물벽 위치 결정을 조건짓는다.

언어모델에는 좌표 원시 데이터 대신 기하량이 미리 계산된 전장 상태를 텍스트로 전달한다.
각 적 무리의 요격점까지의 이동거리와 선회량, 모선 기준 방위, 설치된 그물에 의한 접근로 차단
여부가 포함된다. 출력은 자유 문장이 아니라 다음 자료형이 정하는 JSON
형식으로 강제한다~\citep{openai2026function}. 자유 문장은 파싱에 실패하면 그 주기의 배정을
통째로 잃지만, 이 형식은 디코딩 단계에서 문법으로 작용해 어긋난 토큰 자체를 막는다.

\begin{quote}\ttfamily\scriptsize\setlength{\parindent}{0pt}
CommanderPlan \{\\
\hspace*{1.2em} rationale: str\\
\hspace*{1.2em} deployments: List[\{\\
\hspace*{3.0em} cluster\_id: int,\\
\hspace*{3.0em} ally\_ids: List[int] \}]\\
\hspace*{1.2em} hold\_ships: List[int]\\
\}
\end{quote}

언어모델이 정하는 것은 어느 배가 어느 무리를 맡는가뿐이고 나머지는 대기다. 그물 전개
여부·벽의 방위·교전 반경과 요격점 기하는 기동 계층이 정하며, 배정은 코드가 보정하지 않고
그대로 쓰인다.

한 무리에 두 척 이상이 배정되면 \ref{sec:heuristic}절의 잉여 배정과 같은 반경·방위 분리가
작동하되, 요격점이 요격 구역을 벗어나 유효 마스크가 붕괴하지 않도록 반경을 구역 안으로 자른다.
배정이 바뀌면 요격점과 함께 기동 계층의 유효 마스크(\ref{sec:validmask}절)도 바뀌므로, 배정은
기동의 선행 조건이다.
